\begin{thebibliography}{99}
\raggedright

\bibitem{SourceManuscript}
ProveIt project,
\emph{Polylogarithms manuscript}, canonical modular manuscript,
\texttt{Analysis/Polylogarithms/docs/manuscript},
repository revision
\texttt{f7e35be357654198c1546fd40f363694c5bcc2cc}.
\url{https://github.com/VladimirReshetnikov/ProveIt}.
The source equation \texttt{double:eq:fam2-example} is in
\texttt{chapters/04-depth.tex}; the coincident Hurwitz kernel is in
\texttt{chapters/07-shifted-stieltjes-correlations.tex}.

\bibitem{SourceCanonical}
ProveIt project,
\emph{Canonical Digamma Polynomials and the Fifth Harmonic Bridge},
incoming report in
\path{ProveIt_Canonical_Digamma_Fifth_Moment_2026-10-11.zip},
at the same pinned revision.
Sections 3, 7, and 9 contain the Gamma bridge, exact even-tail
coordinates, and Conjecture \texttt{conj:firstq}.
Archive SHA-256:
\path{3c2fa00434f0366089b819bad77c5a176443594ed1fe6dd0731c877bc74d01a7}.

\bibitem{SourceHigher}
ProveIt project,
\emph{Higher Digamma Powers and Ordered Gamma Sums},
incoming report in
\path{ProveIt_Higher_Digamma_Powers_2026-10-11.zip},
at the same pinned revision.
The all-power algorithm, fifth moment, and normalized primitives are
prior baseline results.

\bibitem{bip:Nakamura2012}
T. Nakamura,
\emph{A simple proof of the functional relation for the Lerch type
Tornheim double zeta function},
Tokyo Journal of Mathematics \textbf{35} (2012), 333--337.
\url{https://arxiv.org/abs/1012.1144}.
See Proposition 1.2, read with the product correction recorded by
Panzer in footnote 16 of the next reference.

\bibitem{bip:Panzer2017}
E. Panzer,
\emph{The parity theorem for multiple polylogarithms},
Journal of Number Theory \textbf{172} (2017), 93--113.
\url{https://arxiv.org/abs/1512.04482}.
See Theorem 1.3, equation (3.2), and footnote 16.

\bibitem{Guo2014}
L. Guo, G. Regensburger, and M. Rosenkranz,
\emph{On integro-differential algebras},
Journal of Pure and Applied Algebra \textbf{218} (2014), 456--473.
\url{https://doi.org/10.1016/j.jpaa.2013.06.015};
\url{https://arxiv.org/abs/1212.0266}.
Proposition 4.10 gives the functional-monomial direct-sum decomposition.

\bibitem{Rosenkranz2008}
M. Rosenkranz and G. Regensburger,
\emph{Integro-differential polynomials and operators},
Proceedings of ISSAC 2008, 261--268.
The reduction rule is given in Lemma 5; the later treatment
\cite{Guo2014} supplies the direct-sum formulation used here.

\bibitem{Shpot2016}
M. A. Shpot, M. P. Chaudhary, and R. B. Paris,
\emph{Integrals of products of Hurwitz zeta functions and the Casimir
effect in $\phi^4$ field theories},
arXiv:1603.00722v2 (2016).
\url{https://arxiv.org/abs/1603.00722}.
Remark 1 states the classical Mikol\'as product integral.  Theorem 2
concerns the shifted, ordinary product integral, which is distinct
from the coordinate finite part evaluated in this article.

\bibitem{DLMF}
NIST Digital Library of Mathematical Functions,
\url{https://dlmf.nist.gov/}.
Relevant formulas and sections: 5.5.4 (digamma reflection),
5.7 (Gamma and digamma series), 5.11.2 and 5.15.9
(digamma and polygamma asymptotics), 24.8.1
(Bernoulli Fourier series), and 25.11 (Hurwitz zeta).
Accessed 11 October 2026.

\end{thebibliography}
